\documentclass{article}
\usepackage{amsmath,amssymb}
\usepackage[vocab={permutations}]{leanify}
\newtheorem{lemma}{Lemma}
\begin{document}

\section{Notation}

We write $S_n$ for the set of permutations of $\{1,\dots,n\}$.
We write $\operatorname{id}_n$ for the identity permutation of $\{1,\dots,n\}$.
We write $\operatorname{fix}(\sigma)$ for the number of fixed points of $\sigma$.
We write $\operatorname{ord}(\sigma)$ for the order of $\sigma$.
We write $\operatorname{inv}(\sigma)$ for the number of inversions of $\sigma$.

\section{Permutations}

\begin{lemma}\label{sign}
For every positive integer $n$ and all $\sigma,\tau\in S_n$, $\operatorname{sgn}(\sigma\tau)=\operatorname{sgn}(\sigma)\operatorname{sgn}(\tau)$.
\end{lemma}

\begin{lemma}\label{inversions}
For every permutation $\sigma$, $\operatorname{sgn}(\sigma)^2=1$ and $\operatorname{inv}(\sigma)\ge0$.
\end{lemma}

\begin{lemma}\label{order}
For every permutation $\sigma$, $\sigma^{\operatorname{ord}(\sigma)}$ is the identity.
\end{lemma}

\begin{lemma}\label{transposition}
Let $\tau$ be a transposition. Then $\tau$ is odd, and $\tau^2$ is the identity.
\end{lemma}

\begin{lemma}\label{fixed}
Let $n$ be a positive integer and let $\sigma$ be a permutation of $\{1,\dots,n\}$. Then $\operatorname{fix}(\sigma)\le n$, and $\sigma(i)\le n$ for every integer $i$ with $1\le i\le n$.
\end{lemma}

\begin{lemma}\label{derangement}
Every derangement $\sigma$ satisfies $\operatorname{fix}(\sigma)=0$.
\end{lemma}

\end{document}
